
\section{Why Deep Learning Still Puzzles Us}

Deep learning has achieved striking empirical success across a wide range of domains:
vision, language, speech, control, scientific modeling, and generative AI.
Yet this success raises a scientific puzzle.

\medskip

Classical learning intuition suggests that highly flexible models should be at serious risk of overfitting.
Deep neural networks are often massively overparameterized, capable of fitting training data extremely well, and sometimes even capable of interpolating the training sample almost perfectly.
Nevertheless, they frequently generalize far better than naive intuition would predict.

\medskip

This creates a tension:

\begin{quote}
deep learning works extremely well in practice, but our theory still does not provide a single complete explanation of why it works so well.
\end{quote}

\subsection*{Overparameterization and interpolation}

A modern neural network may contain far more parameters than training examples.
From a purely classical counting perspective, this might appear dangerous.
One could imagine many parameter settings that fit the training data while behaving very differently away from it.

\medskip

And yet empirical practice shows something subtler.
In many regimes, larger networks are not only more expressive, but also easier to optimize and sometimes better at generalizing.

\medskip

This does not mean classical learning theory is wrong.
Rather, it means that the interaction between model size, optimization, data structure, and inductive bias is more subtle than early intuition suggested.

\subsection*{Optimization and generalization are different questions}

Two questions must be kept separate:
\begin{enumerate}
    \item \textbf{Optimization:} can we find parameters that fit the training objective well?
    \item \textbf{Generalization:} will the resulting solution perform well on new data?
\end{enumerate}

Deep learning has complicated both questions.
Overparameterization may make optimization easier in some regimes, while also changing the geometry of the space of solutions.
The learned model is not just any interpolating solution; it is the solution found by a particular optimization process with a particular architecture and initialization.

\subsection*{Why theory trails practice}

One reason theory trails practice is that modern systems combine many ingredients:
\begin{itemize}
    \item compositional function classes,
    \item stochastic gradient-based optimization,
    \item normalization and residual pathways,
    \item architectural inductive bias,
    \item large-scale pretraining,
    \item self-supervised and generative objectives.
\end{itemize}

Each ingredient is partly understood in isolation.
What is harder is to explain how they interact in realistic large systems.

\subsection*{The purpose of this chapter}

This chapter is not meant to give a complete theory of deep learning.
That theory does not yet exist in a fully unified form.
Instead, the goal is more modest and more useful:
\begin{itemize}
    \item to identify what is genuinely puzzling,
    \item to describe the main partial explanations,
    \item to explain why strong benchmark performance is not the same as reliability,
    \item to clarify what remains open.
\end{itemize}

